% manuscript.tex -- Journal of Systems and Software submission.
% Special Issue "AI Techniques for Performance, Reliability, and
% Sustainability of Modern Software Systems" (VSI:AI4MSS).

\documentclass[preprint,3p]{elsarticle}

\usepackage[T1]{fontenc}
\usepackage{lmodern}
\usepackage{array}
\usepackage{calc}
\usepackage{amsmath,amssymb}
\usepackage{graphicx}
\usepackage{booktabs}
\usepackage{multirow}
\usepackage{longtable}
\usepackage{lineno}
\usepackage{xcolor}
\usepackage{textcomp}
\usepackage{xr-hyper}
\usepackage{hyperref}
% Supplement labels are referenced as S-<label>; this reads supplementary.aux,
% so build the supplement before the final pass (see Makefile).
\externaldocument[S-]{supplementary}
\usepackage{microtype}

% Float packing: the default Elsevier preprint parameters strand large tables and
% leave partially-filled pages. These are typesetting settings only.
\renewcommand{\topfraction}{0.9}
\renewcommand{\bottomfraction}{0.8}
\renewcommand{\textfraction}{0.07}
\renewcommand{\floatpagefraction}{0.75}
\setcounter{topnumber}{3}
\setcounter{bottomnumber}{2}
\setcounter{totalnumber}{5}

\setlength{\textfloatsep}{7pt plus 2pt minus 2pt}
\setlength{\intextsep}{7pt plus 2pt minus 2pt}
\setlength{\floatsep}{6pt plus 2pt minus 2pt}
\setlength{\abovecaptionskip}{3pt}
\setlength{\belowcaptionskip}{1pt}

% Experiment pages in the public replication repository, pinned to the submission tag.
\newcommand{\sagexperimentsurl}{https://github.com/onuralpyigit/software-as-a-graph/tree/jss-submission-v4/docs/research/jss/experiments}
\newcommand{\sagexperiments}{\url{\sagexperimentsurl}}

\providecommand{\tightlist}{%
  \setlength{\itemsep}{0pt}\setlength{\parskip}{0pt}}

\journal{Journal of Systems and Software}

\begin{document}

\begin{frontmatter}

\title{Software-as-a-Graph: Predicting Cascading-Failure Impact in
  Publish--Subscribe Systems Before Deployment with QoS-Aware Graphs
  and Hybrid Learning}

\author[itu]{Ibrahim Onuralp Yigit\corref{cor1}}
\ead{yigiti@itu.edu.tr}
\author[itu]{Feza Buzluca}
\ead{buzluca@itu.edu.tr}

\affiliation[itu]{organization={Department of Computer Engineering, Istanbul Technical University},
             city={Istanbul}, postcode={34469}, country={Turkey}}
\cortext[cor1]{Corresponding author}

\begin{abstract}
Publish--subscribe middleware decouples components in space and time, and in doing so hides the paths along which failures cascade. Architects need to know which components are systemically critical before deployment, when no runtime telemetry exists. We present Software-as-a-Graph (SaG), a framework that turns architecture descriptions into typed multigraphs with Quality-of-Service (QoS) weighted dependencies, ranks components by cascading-failure impact, and explains each flagged component with an ISO/IEC 25010 reliability and maintainability profile. SaG provides a training-free closed-form engine, graph-learning engines with 16-dimensional QoS edge encodings, and hybrid engines in which a learned model corrects the closed-form score. Under leave-one-scenario-out cross-validation over twelve synthetic architectures, the QoS-aware representation alone raises closed-form ranking from Spearman $\rho = 0.349$ to $0.553$, on every held-out architecture. Closed-form and learned engines are strongest on different architectures, and neither significantly beats the other alone. The hybrids that combine them are the most accurate engines on held-out architectures ($\rho = 0.657$ and $0.683$) and outperform the closed-form engine on 11 of 12 folds, under decision rules registered before their runs. Trained only on synthetic data, learned engines transfer zero-shot to hand-authored models of five open-source systems at $\rho = 0.760$--$0.805$, against $0.511$--$0.526$ for every training-free score, and roughly double top-$K$ critical-set overlap. A capacity-matched control shows that the QoS edge encoding, not relation-specific parameters, drives learned accuracy. Neural inference takes milliseconds. Declaring QoS contracts in architecture models thus makes cascading-failure risk measurable, and explainable, before deployment.

\end{abstract}

\begin{keyword}
Graph neural networks \sep
dependability \sep
publish--subscribe \sep
cascading failures \sep
hybrid learning \sep
quality of service \sep
explainable AI
\end{keyword}

\end{frontmatter}

\linenumbers

\section{Introduction}
\label{sec:1}

\subsection{Motivation and Problem}
\label{sec:1.1}

Large distributed systems increasingly communicate through asynchronous publish--subscribe (pub-sub) middleware: ROS~2 in autonomous driving \cite{macenski2022robot}, Apache Kafka in enterprise event streams \cite{kreps2011kafka}, DDS in cyber-physical systems \cite{omg2015dds}, MQTT in IoT \cite{oasis2019mqtt}, and cloud-native microservices \cite{dragoni2017microservices,newman2015building}. Pub-sub decouples producers and consumers in space, time and synchronization \cite{eugster2003faces}. Components interact through topics and brokers rather than direct references, and deployment-time Quality-of-Service (QoS) policies govern reliability, durability, priority and deadlines.

The same decoupling hides how failures spread. Publishers and subscribers share no direct link, so outages, head-of-line blocking and backpressure propagate along concealed paths through brokers, shared topics, colocated hosts and shared libraries \cite{motter2002cascade,buldyrev2010cascade}. These failures take two forms. In \emph{sequential cascades}, a slow subscriber fills a broker queue and gradually starves its publishers \cite{albert2000error}. In \emph{simultaneous blasts}, a shared-library crash or host outage takes down every colocated service at once. Neither architecture diagrams nor static call graphs show these mechanisms. The cheapest time to reduce the risk is before deployment, at design and continuous-integration time \cite{avizienis2004basic,bass2012software}, when no runtime telemetry exists. Architects therefore need to know, from configuration manifests alone, which components, topics and links are systemically critical and why.

Existing practice leaves this gap open, which we call the \textbf{Architecture--Code Gap}: a system can have bug-free code in every service and still be fragile through hidden single points of failure or mismatched QoS contracts \cite{perry1992foundations,garcia2009identifying}. Architecture evaluations such as ATAM rely on manual elicitation \cite{kazman1998atam}. Static code analysis inspects services in isolation \cite{sonarsource2024clean,chidamber1994metrics}. Chaos engineering needs a provisioned cluster \cite{basiri2016chaos}. Homogeneous centrality flattens typed topologies into untyped graphs \cite{freeman1977centrality,brandes2001faster}. Learned models could combine these structural cues, but on their own they produce risk scores without actionable explanations. What is missing is a representation that makes pub-sub failure paths explicit, and evidence on which analyzer to trust on it.

\subsection{The Software-as-a-Graph (SaG) Approach}
\label{sec:1.3}

\textbf{Software-as-a-Graph (SaG)} is a pre-deployment static system analysis framework for event-driven architectures (Figure~\ref{fig:1}). It (1)~models an architecture as a typed, directed multigraph over Applications, Brokers, Topics, Execution Hosts and Shared Libraries (Section~\ref{sec:3.1}); (2)~derives a QoS-weighted \texttt{DEPENDS\_ON} projection that captures both sequential cascades and simultaneous blasts (Section~\ref{sec:3.2}); (3)~ranks components by predicted cascade impact with a training-free closed-form engine, graph-learning engines, and hybrid engines in which a learned model corrects the closed-form score (Section~\ref{sec:4}); and (4)~explains flagged components with an ISO/IEC 25010 Reliability--Maintainability (RM) profile (Section~\ref{sec:5}). Predictors read only the analysis graph. Ground-truth impact comes from independent simulators that run on the raw structural topology (Section~\ref{sec:4.4}). The central thesis is that the representation, not the choice of analyzer, is the decisive investment: once declared QoS contracts are part of the graph, closed-form and learned engines both rank failure impact better, and each engine is the best choice in a distinct deployment context.

\subsection{Research Questions}
\label{sec:1.4}

\begin{itemize}\tightlist
\item \textbf{RQ1 (Ranking accuracy):} \emph{How accurately do SaG's closed-form, learned and hybrid engines rank components by cascading-failure impact on unseen architectures, compared with standard structural baselines?}
\item \textbf{RQ2 (What learning needs):} \emph{Does the learned engines' accuracy come from relation-specific (typed) parameters or from the QoS edge encoding, once model capacity and edge-channel width are matched?}
\item \textbf{RQ3 (Transfer):} \emph{How well do engines trained on synthetic architectures transfer zero-shot to independently authored models of five open-source systems?}
\item \textbf{RQ4 (Cost):} \emph{What does the analysis cost at CI/CD time, which stage dominates, and how does it compare with running the simulation directly?}
\end{itemize}

\noindent The primary contrast, the matched control and both hybrid engines were each registered with a decision rule before their results existed. Supplementary Section~\ref{S-supp:amendments} logs every later change as an amendment and maps these four questions onto the five of the registered analysis plan.

\subsection{Contributions}
\label{sec:1.5}

Evaluated under leave-one-scenario-out (LOSO) cross-validation over twelve synthetic architectures and zero-shot on five open-source system models, this paper contributes:

\begin{enumerate}\tightlist
    \item \textbf{A QoS-aware typed architecture model} that derives logical dependencies from physical pub-sub linkages, weights them by declared QoS contracts, and distinguishes sequential cascades from simultaneous blasts (Section~\ref{sec:3}). This representation is the largest single gain in the study: ranking on it raises Spearman correlation with simulated cascade impact from $0.349$ (unweighted centrality) to $0.553$, on all twelve held-out architectures ($+0.204$, $p = 0.0005$).
    \item \textbf{Complementary closed-form, learned and hybrid engines} (Section~\ref{sec:4}). Closed-form and learned engines are strongest on different architectures, and neither significantly outperforms the other alone. Hybrid engines that learn a correction to the closed-form score exploit this complementarity. They are the most accurate engines on held-out architectures ($\rho = 0.657$ and $0.683$; $+0.103$ and $+0.130$ on 11 of 12 folds). Each meets the decision rule registered before its run and remains significant under a Holm correction pooled over all eleven registered contrasts.
    \item \textbf{Evidence on what graph learning needs and how far it transfers} (Section~\ref{sec:7}). A capacity- and channel-matched control shows that the 16-D QoS edge encoding, not relation-specific weights, drives learned accuracy ($+0.073$ on 10 of 12 folds) and stabilizes training, so an untyped attention network suffices. Trained only on synthetic data, learned engines transfer zero-shot to five open-source system models at $\rho = 0.760$--$0.805$, against $0.511$--$0.526$ for every training-free score, and nearly double top-$K$ critical-set overlap.
    \item \textbf{A standards-grounded explanation layer} (Section~\ref{sec:5}) that attributes each flagged component to ISO/IEC 25010 Availability, Fault Tolerance or Maintainability, and so names the remediation it calls for: replication, circuit breakers or decoupling.
    \item \textbf{A reproducible benchmark and cost profile}: seventeen architectures totaling 2,812 components, twelve of which regenerate byte-identically from committed configurations. Every reported table value is mechanically reconciled against released artifacts. Neural inference takes $56\,\text{ms}$ on a 2,000-component architecture, and one structural metric dominates cost (Section~\ref{sec:rq4}).
\end{enumerate}

\noindent A previous conference paper \cite{yigit2025graph} introduced the preliminary multigraph and deterministic quality model on synthetic topologies. This paper adds the learned and hybrid engines, the QoS edge encoding, LOSO and zero-shot evaluation, the matched control, and the cost profile, and it repositions that quality model as the explanation layer.

Section~\ref{sec:2} reviews related work, Sections~\ref{sec:3}--\ref{sec:5} present the model, engines and explanation layer, Sections~\ref{sec:6}--\ref{sec:7} the evaluation, Section~\ref{sec:8} the discussion and threats, and Section~\ref{sec:9} concludes.

\section{Related Work}
\label{sec:2}

\subsection{Dependability Analysis of Distributed Systems}
\label{sec:2.1}

Runtime approaches to dependability, such as broker clustering, backpressure, autoscaling, failover and chaos engineering \cite{basiri2016chaos}, require a running cluster, can disrupt service, and consume substantial compute. That compute is itself a concern of green software engineering \cite{calero2015green,lannelongue2021green,verdecchia2023greenai}. Architecture-based reliability prediction has a long history. Cheung's absorbing Markov chain \cite{cheung1980user} and the state-, path- and additive models surveyed by Goseva-Popstojanova and Trivedi \cite{gosevapopstojanova2001architecture} and Immonen and Niemel{\"a} \cite{immonen2008survey} are examples, as are the Palladio Component Model \cite{becker2009palladio}, layered queueing networks \cite{franks2009layered} and the AADL Error Model Annex \cite{delange2014aadl}. These approaches answer broader questions but need operational profiles and failure rates that are unavailable at commit time. SaG asks a narrower question from manifests alone: whose failure reaches furthest in the declared topology?

Telemetry-driven methods diagnose faults in running microservices. Examples are Seer \cite{gan2019seer} and Sage \cite{gan2021sage}, MicroRCA \cite{wu2020microrca}, TraceRCA \cite{li2021tracerca}, and GNN-based DeepTraLog \cite{zhang2022deeptralog}, Eadro \cite{lee2023eadro} and MicroCause \cite{meng2020microcause} (reviewed in \cite{zhang2024failurediag}). In synchronous microservices, cascades stem from thread-pool exhaustion, RPC timeouts and retry storms \cite{zhou2021fault}. In pub-sub systems they spread through queue saturation and message starvation. All of these methods need traces or metrics from a live system; SaG works before any code runs. SaG's question is also that of change impact analysis \cite{bohner1996impact} and of dependency management in microservice fleets \cite{esparrachiari2018tracking}. Ranking by impact relates to learning-to-rank defect prediction \cite{yang2015ltr}, which optimizes the ranking measure directly, as our listwise loss does (Section~\ref{sec:4.2}).

\subsection{Static Code Analysis and Static System Analysis}
\label{sec:2.2}

Static code analysis (SCA) tools such as SonarQube \cite{sonarsource2024clean} measure complexity \cite{mccabe1976complexity}, cohesion and coupling \cite{chidamber1994metrics,fenton2014metrics} within individual services to flag defect-prone modules \cite{basili1996validation,nagappan2005static,zimmermann2007predicting,menzies2007data}. SCA cannot see inter-service messaging, broker saturation or cross-host propagation. Architecture recovery tools reconstruct system-level structure from code \cite{bushong2021static,schneider2024comparison}, and HAROS \cite{santos2019haros} checks ROS systems statically before launch. SaG's static system analysis (SSA) uses the declared topology instead, propagating code-level metrics across architectural dependencies. This lets teams find structural anti-patterns \cite{garcia2009catalogue,taibi2018microservice} and architectural technical debt \cite{li2015technicaldebt} in CI/CD \cite{humble2010continuous,chen2015continuous}.

\subsection{Quality Models and Multi-Criteria Evaluation}
\label{sec:2.3}

ISO/IEC 25010:2023 \cite{iso25010_2023} and ISO/IEC 25019:2023 \cite{iso25019_2023} define product quality and quality in use. SaG covers the characteristics derivable from deployment topology, namely Availability, Fault Tolerance and the Maintainability sub-characteristics (Section~\ref{sec:5.1}), and links internal structural quality to external dependability \cite{iso25023_2016,iso25021_2012}. Aggregating metrics into an auditable score is a multi-criteria decision problem, for which the Analytic Hierarchy Process (AHP) \cite{saaty1980ahp} is standard. AHP's consistency ratio detects inconsistent judgments but not matrices back-filled from a chosen answer, a distinction this study reports for its weights (Supplementary Section~\ref{S-supp:ahp}).

\subsection{Graph Learning and Explainability}
\label{sec:2.4}

Centrality indices \cite{freeman1977centrality,brandes2001faster,brin1998pagerank,newman2010networks} and cascade models of network robustness \cite{albert2000error,motter2002cascade,buldyrev2010cascade} assume homogeneous, usually undirected graphs. A single untyped score conflates structurally different elements, such as topics, libraries and hosts, and cannot say \emph{why} a component is critical. Learned node-importance methods, including FINDER \cite{fan2020keyplayers}, DrBC \cite{fan2019learning} and PowerGraph \cite{varbella2024powergraph}, share the homogeneity assumption. Homogeneous GNNs (GCN \cite{kipf2017semi}, GraphSAGE \cite{hamilton2017inductive}, GAT \cite{velickovic2018graph}) discard relation identity unless it is supplied as a feature. Heterogeneous GNNs (RGCN \cite{schlichtkrull2018rgcn}, HAN \cite{wang2019han}, HGT \cite{hu2020hgt}, MAGNN \cite{fu2020magnn}) learn relation-specific transformations. We evaluate both families under matched capacity (Section~\ref{sec:rq2}). Khodabandeh et al.~\cite{khodabandeh2025linkpred} apply graph attention to microservice call graphs to predict future interactions. We instead predict the impact of removing a node from a declared topology. GNN explainers such as GNNExplainer \cite{ying2019gnnexplainer} and PGExplainer \cite{luo2020pgexplainer} explain models in terms of their internal features. SaG's explanation layer (Section~\ref{sec:5}) instead names ISO/IEC quality sub-characteristics and a remediation class for each flagged component.

\section{The Software-as-a-Graph (SaG) Architectural Model}
\label{sec:3}

Figure~\ref{fig:1} presents the end-to-end architecture of the SaG framework. The shared front end processes Architecture-as-Code manifests, constructs a typed multigraph, projects implicit runtime interactions throughout a QoS-weighted logical dependency layer, and extracts typed node properties. These features feed the ranking engines (Section~\ref{sec:4}) and, separately and with no shared parameters, the explanation layer (Section~\ref{sec:5}).

\begin{figure}[htbp]

\centering

\includegraphics[width=0.70\linewidth]{figures/Figure_1}

\caption{End-to-end architecture of the SaG framework. The predictive pathway runs down the centre: manifest ingestion, typed multigraph, QoS-weighted \texttt{DEPENDS\_ON} projection with typed node properties, the ranking engines (closed-form, learned and hybrid; Figure~\ref{fig:engines}), and the ranked critical set. The dashed edge marks the ground-truth simulation oracle, which operates only on $G_{\text{structural}}$, trains the predictor offline and takes no part in inference. The explanation layer re-enters from the analysis multigraph and shares no parameters with the predictor, reaching flagged components through triage rather than data flow.}

\label{fig:1}

\end{figure}

\subsection{Formal Multigraph Definition}
\label{sec:3.1}

A distributed system is described as a typed, weighted, directed multigraph:

\begin{equation}
\label{eq:sag_multigraph}
\mathcal{G} = (V, E, \tau_V, \tau_E, w_V, w_E)
\end{equation}

where:

\begin{itemize}\tightlist
    \item $V = V_{\text{app}} \cup V_{\text{broker}} \cup V_{\text{topic}} \cup V_{\text{host}} \cup V_{\text{lib}}$ holds the five entity types $\mathcal{T}_V$ of Table~\ref{tab:1}; $V_{\text{host}}$ denotes physical or virtual \emph{Execution Hosts}.
    \item $E$ is the set of directed edges, and $\tau_V: V \to \mathcal{T}_V$ and $\tau_E: E \to \mathcal{T}_E$ assign entity and relation types.
    \item $w_V: V \to (0, 1]$ and $w_E: E \to (0, 1]$ weight entity criticality and connection strength. For Applications and Libraries, $w_V(v) = 1 - \text{CQP}(v)$, where CQP is a code-quality penalty computed from static code metrics (lines of code, cyclomatic complexity, coupling and cohesion); otherwise $w_V(v) = 1.0$.
\end{itemize}

\begin{table}[htbp]
\centering
\small
\caption{Entity types and structural edge types in the SaG model.}
\label{tab:1}
\resizebox{\linewidth}{!}{%
\begin{tabular}{lll}
\toprule
\textbf{Entity Type ($\mathcal{T}_V$)} & \textbf{Architectural Role} & \textbf{Concrete System Examples} \\
\midrule
\textbf{Application} ($V_{\text{app}}$) & Process producing/consuming messages & ROS~2 node, Kafka microservice, MQTT client \\
\textbf{Broker} ($V_{\text{broker}}$) & Message routing and queuing intermediary & RabbitMQ exchange, Mosquitto, EMQX broker \\
\textbf{Topic} ($V_{\text{topic}}$) & Named logical communication channel & \texttt{/sensor/lidar}, \texttt{orders.payment.completed} \\
\textbf{Execution Host} ($V_{\text{host}}$) & Physical or virtualized execution environment & Bare-metal server, Kubernetes worker, Cloud VM \\
\textbf{Library} ($V_{\text{lib}}$) & Shared software package or runtime dependency & Kafka client, OpenCV, Protobuf runtime \\
\midrule
\textbf{Structural Edge ($\mathcal{T}_E$)} & \textbf{Direction} & \textbf{Semantic Meaning} \\
\midrule
\texttt{PUBLISHES\_TO} / \texttt{SUBSCRIBES\_TO} & App/Library $\to$ Topic & Component publishes to / consumes from topic \\
\texttt{ROUTES} & Broker $\to$ Topic & Broker manages and routes topic traffic \\
\texttt{RUNS\_ON} & App/Broker $\to$ Host & Process is hosted on host \\
\texttt{CONNECTS\_TO} & Host $\to$ Host & Network link between hosts \\
\texttt{USES} & App $\to$ Library & Application links to shared library \\
\bottomrule
\end{tabular}%
}
\end{table}

\begin{table}[htbp]
\centering
\small
\caption{Notation used throughout. Entity and edge types: Table~\ref{tab:1}; simulation oracles: Table~\ref{tab:oracle_roles}.}
\label{tab:notation}
\begin{tabular}{@{}l@{\hspace{5pt}}l@{\hspace{10pt}}l@{\hspace{5pt}}l@{}}
\toprule
$G_{\text{structural}}$ & Raw multigraph; oracles only & $Q(v)$ & RM composite quality score \\
$G_{\text{analysis}}$ & \texttt{DEPENDS\_ON} projection; predictor input & $\rho$ & Spearman $\rho$, full population \\
$V_{\text{app}}$ & Application nodes; the scored population & $\rho_{>0}$ & Spearman $\rho$, active stratum ($I^* > 0$) \\
$w(t)$, $w(e)$ & QoS topic weight, edge weight & Overlap@$K$ & Top-$K$ set overlap, $K = 0.20\,|V_{\text{app}}|$ \\
$I^*(v)$ & Primary cascade-reachability oracle & $I_{\text{comp}}$, $I_{\text{dyn}}$, $I_M$ & Further oracles (Table~\ref{tab:oracle_roles}) \\
\bottomrule
\end{tabular}
\end{table}


\subsection{QoS-Aware Weights and Logical Dependency Derivation}
\label{sec:3.2}

A link's strength depends on its QoS contract: a \texttt{RELIABLE} topic with \texttt{TRANSIENT\_LOCAL} durability couples services more strongly than a \texttt{BEST\_EFFORT} telemetry stream. Each topic $t$ therefore carries a weight $w(t) \in (0, 1]$ combining its declared QoS with payload size and publication frequency:

\begin{equation}
\label{eq:3}
w(t) = \alpha_{\text{top}} \cdot \text{QoS}(t) + \beta_{\text{top}} \cdot \text{SizeNorm}(t) + \gamma_{\text{top}} \cdot \text{FreqNorm}(t),
\quad (\alpha_{\text{top}},\, \beta_{\text{top}},\, \gamma_{\text{top}}) = (0.75,\, 0.15,\, 0.10)
\end{equation}
where the QoS term is an AHP-weighted aggregate of the declared contract:

\begin{equation}
\label{eq:topic_qos}
\text{QoS}(t) = w_{\text{rel}} \cdot q_{\text{rel}} + w_{\text{dur}} \cdot q_{\text{dur}} + w_{\text{prio}} \cdot q_{\text{prio}},
\quad (w_{\text{rel}}, w_{\text{dur}}, w_{\text{prio}}) = (0.24,\, 0.62,\, 0.14)
\end{equation}

Here $q_{\text{rel}} \in \{0, 1\}$ (best-effort, reliable), $q_{\text{dur}} \in \{0, 0.5, 1\}$ (volatile, transient-local, persistent) and $q_{\text{prio}} \in \{0, 0.5, 1\}$ (low, medium, high). Durability carries the highest weight because it determines whether state and message history survive restarts. The sub-weights come from a Saaty pairwise matrix with genuine second-eigenvalue spread and $CR = 0.016$ (Supplementary Section~\ref{S-supp:ahp}). Size and frequency are log-compressed and clamped:

\begin{equation}
\label{eq:sizenorm_freqnorm}
\text{SizeNorm}(t) = \min\left(1.0, \frac{\log_2(1 + B(t))}{20}\right), \quad
\text{FreqNorm}(t) = \min\left(1.0, \frac{\log_{10}(1 + F(t))}{3}\right)
\end{equation}

where $B(t)$ is the payload in bytes (design envelope 1~MiB, the practical DDS sample limit) and $F(t)$ the publication frequency in Hz. $w(t)$ is clamped to $[0.01, 1]$ so that best-effort edges remain visible, and every \texttt{PUBLISHES\_TO}, \texttt{SUBSCRIBES\_TO} and \texttt{ROUTES} edge of $t$ carries $w_E(e) = w(t)$ and the topic's QoS vector. The $(\alpha_{\text{top}}, \beta_{\text{top}}, \gamma_{\text{top}})$ split is a documented choice rather than a sensitive parameter: no point of its simplex changes any reported ordering (Supplementary Section~\ref{S-supp:params}).

\subsubsection{Logical Dependency Projection (\texttt{DEPENDS\_ON})}

Structural edges do not capture implicit runtime dependencies: a subscriber depends on a publisher, yet no edge joins them. SaG therefore derives one semantic relation, \texttt{DEPENDS\_ON}, directed from \emph{dependent} to \emph{dependency} (``if the target fails, the source is impacted''), by the six rules of Table~\ref{tab:2}. Its weight $w \in (0, 1]$ expresses how likely a disruption of the dependency is to reach the dependent.

\begin{table}[htbp]
\centering
\small
\caption{The six \texttt{DEPENDS\_ON} logical dependency projection rules.}
\label{tab:2}
\resizebox{\linewidth}{!}{%
\begin{tabular}{clll}
\toprule
\textbf{Rule} & \textbf{Dependency Category} & \textbf{Structural Pattern ($\text{Dependent} \to \text{Dependency}$)} & \textbf{Derived Weight ($w$)} \\
\midrule
\textbf{1} & \texttt{app\_to\_app} & Subscriber $\to$ Publisher (via shared Topic, incl.\ transitive \texttt{USES}) & $1 - \prod_{t \in T}(1 - w(t))$ \\
\textbf{2} & \texttt{app\_to\_broker} & Publisher/Subscriber $\to$ Broker routing its topics & $1 - \prod_{t \in T}(1 - w(t))$ \\
\textbf{3} & \texttt{host\_to\_host} & Host $\to$ Host (lifted from inter-host app dependencies) & $\max_{u \in \text{hosted}(h_1), v \in \text{hosted}(h_2)} w_{\text{DEPENDS\_ON}}(u \to v)$ \\
\textbf{4} & \texttt{host\_to\_broker} & Host $\to$ Broker (lifted from hosted app dependencies) & $\max_{u \in \text{hosted}(h)} w_{\text{DEPENDS\_ON}}(u \to b)$ \\
\textbf{5} & \texttt{app\_to\_lib} & Application $\to$ Shared Library it \texttt{USES} & $H(w_V(\text{app}), w_V(\text{lib}))$ \\
\textbf{6} & \texttt{broker\_to\_broker} & Broker $\leftrightarrow$ Broker (shared physical fault-domain colocation, symmetric) & $w_V(\text{host})$ \\
\bottomrule
\end{tabular}%
}
\end{table}

Rules 1 and 2 combine the topics $T$ joining a pair by probabilistic union rather than maximum \cite{pearl1988probabilistic,beliakov2007aggregation,yager1988owa}, so parallel failure paths always increase coupling. Rule 5 uses the harmonic mean $H(x, y) = 2xy/(x+y)$ \cite{hardy1952inequalities}, and Rules 3 and 4 lift dependencies to hosts by maximum.

\noindent\textbf{Sequential cascades and simultaneous blasts.} Rule~1 captures sequential cascades, in which a failed publisher starves subscribers through queues and topic buffers. Rule~5 captures simultaneous blasts, in which a crashed library or host takes down every consumer at once. Untyped graphs collapse the two into indistinguishable edges. Rule~6, the only symmetric rule, joins brokers colocated on a host, which share its failure domain. Figure~\ref{fig:2} shows both mechanisms on a seven-entity example.

\begin{figure}[htbp]
\centering
\includegraphics[width=\linewidth]{figures/Figure_2}
\caption{Running example. (a)~Three applications share topic $t$ (routed by broker $b$) and library $\ell$, and all run on host $n$. No structural edge joins two applications. (b)~The derived \texttt{DEPENDS\_ON} edges make the hidden dependencies explicit: the subscribers $a_2, a_3$ depend on the publisher $a_1$ (Rule~1, a sequential cascade through the topic), every application depends on $\ell$ (Rule~5, a simultaneous blast if $\ell$ fails), and each application depends on the broker routing its topic (Rule~2). Simulation oracles run on view~(a) only; predictors read view~(b).}
\label{fig:2}
\end{figure}

\subsection{Dual Graph Views}
\label{sec:3.3}

The \textbf{structural graph} $G_{\text{structural}}$ is the raw deployment topology. The \textbf{analysis graph} $G_{\text{analysis}}$ adds the derived, QoS-weighted \texttt{DEPENDS\_ON} edges and the code metrics (Figure~\ref{fig:2}). All predictor features are computed on $G_{\text{analysis}}$, while simulation oracles run only on $G_{\text{structural}}$ (Section~\ref{sec:4.4}); Rule~6 therefore cannot influence any label.

\subsection{Typed Node Feature Encoding}
\label{sec:3.4}

Both the predictive pathway (Section~\ref{sec:4}) and the explanation layer (Section~\ref{sec:5}) read the same typed node properties from $G_{\text{analysis}}$: the predictor projects them per entity type before message passing, the explanation layer aggregates them into a quality profile. All five entity types share a deterministic 18-dimensional base block of topological metrics, each normalized to $[0, 1]$ within its graph so that raw graph size does not drive cross-scenario transfer. The block comprises PageRank and Reverse PageRank; betweenness, closeness and eigenvector centrality; in- and out-degree; clustering; articulation and bridge scores; the node QoS weight and QoS-weighted in- and out-degree; multi-path coupling; path complexity; fan-out criticality; and the Connectivity Degradation Index (CDI), which removes each node and measures the resulting connectivity loss (full schema: Supplementary Section~\ref{S-supp:features}). Type-specific blocks extend the vector to 19--25 dimensions: code metrics and CQP for Applications, reverse-\texttt{USES} blast radius for Libraries, queue capacity for Brokers, publisher/subscriber counts and QoS criticality for Topics, and CPU and memory for Hosts. PageRank, Reverse PageRank, betweenness and eigenvector centrality are computed on the QoS-weighted projection. They therefore carry QoS information into every predictor, including the ``QoS-off'' arms, which lack only the explicit QoS edge channel and QoS node columns. Because these topological summaries are available to any scorer, closed-form engines are genuine competitors rather than strawmen. CDI, at $O(|V|^2 + |V||E|)$, dominates analysis cost (Section~\ref{sec:rq4}).

\section{Ranking Engines and Ground Truth}
\label{sec:4}

SaG ranks components with three kinds of engine. The \textbf{closed-form engine} \texttt{Topo-QoS} is QoS-weighted betweenness on the dependency projection (Section~\ref{sec:6.2}). The \textbf{learned engines} are graph neural networks over the typed multigraph (this section). The \textbf{hybrid engines} are learned engines that correct the closed-form score (Section~\ref{sec:rq1-hybrid}). All are trained or scored against simulation oracles that run on a separate graph view (Sections~\ref{sec:4.3}--\ref{sec:4.4}). Figure~\ref{fig:engines} shows how the three engines relate and how they are evaluated. Full hyperparameters and training commands are on the experiment pages of the replication repository (Section~\ref{sec:6.1}).

\begin{figure}[htbp]
\centering
\includegraphics[width=\linewidth]{figures/Figure_3}
\caption{(a)~SaG's three ranking engines read the same analysis graph. The closed-form engine scores QoS-weighted betweenness $p(v)$; the learned engine outputs a logit $z(v)$. A hybrid engine gives the learned engine $p(v)$ as an extra input feature and adds a learned correction to it on the logit scale, $\sigma(z + \alpha\,\operatorname{logit} p)$, with one learnable scalar $\alpha$. (b)~Ground truth comes from simulation oracles on the structural graph, which no predictor reads. Engines are evaluated by leave-one-scenario-out cross-validation over twelve synthetic architectures (each row trains on eleven and tests on the held-out one) and zero-shot on five open-source system models.}
\label{fig:engines}
\end{figure}

\subsection{Heterogeneous Graph Transformer}
\label{sec:4.1}

The learned engine \texttt{HGT-QoS} is a three-layer Heterogeneous Graph Transformer (HGT) \cite{hu2020hgt} in PyTorch Geometric \cite{fey2019fast}, with hidden dimension $D = 64$ and $H = 4$ heads. Entity-specific projections map raw features $x_v \in \mathbb{R}^{19\text{--}25}$ (Section~\ref{sec:3.4}) into the hidden space, $h_v^{(0)} = \text{LayerNorm}(\text{GELU}(W_{\tau(v)} x_v))$. For each meta-relation $\langle \tau(u), \phi(e), \tau(v)\rangle$, attention uses type-parameterized keys $K(u) = h_u W_K^{\tau(u)}$, queries $Q(v) = \tilde{h}_v W_Q^{\tau(v)}$ and values $V(u) = h_u W_V^{\tau(u)}$, scaled by a learned per-meta-relation prior $\mu_{\langle \tau(u), \phi(e), \tau(v)\rangle}$. That prior is the relation-typing mechanism that RQ2 tests. Message passing runs over both $G_{\text{analysis}}$ and its transpose, to capture downstream starvation and upstream backpressure, with residual connections, dropout $0.10$ and layer normalization.

\subsubsection{QoS Edge Encoding (16-D)}
\label{sec:4.1.1}

Each directed edge carries a vector $e_{uv} \in \mathbb{R}^{16}$. Index 0 is the coupling weight $w_E(e)$ (Section~\ref{sec:3.2}), index 1 the normalized count of simple paths through $e$, indices 2--8 a one-hot of the seven relation types, and indices 9--15 the middleware QoS parameters on \texttt{PUBLISHES\_TO}/\texttt{SUBSCRIBES\_TO} edges, zero elsewhere. Six QoS dimensions are active in our corpus: reliability, durability, priority, a flag for edges whose QoS departs from the scenario's modal profile, and a deadline pair (active flag and log-deadline, populated on $463$ of $615$ topics). A seventh, max-blocking time, is reserved for hard real-time profiles and is zero throughout. The encoding is projected and added to the target representation before attention, $\tilde{h}_v = h_v + W_{\text{edge}} e_{uv}$.

\subsection{Prediction Head and Training Objective}
\label{sec:4.2}

A composite head predicts cascade impact, $\hat{I}^*(v) = \sigma(\text{MLP}_C(h_v \parallel \hat{a}_1(v) \parallel \hat{a}_2(v)))$. The two auxiliary heads $\hat{a}_1, \hat{a}_2$ act only as learned feature enrichment: $\hat{a}_1$ is supervised on $I^*$ and $\hat{a}_2$ is unsupervised. The optimized objective combines regression with listwise and pairwise ranking:
\begin{equation}
\label{eq:loss_active}
\mathcal{L} = \text{MSE}(\hat{I}^*, I^*) + 0.5 \cdot \text{MSE}(\hat{a}_1, I^*) + 0.3 \cdot \mathcal{L}_{\text{rank}} + 0.1 \cdot \mathcal{L}_{\text{pairwise}},
\end{equation}
where $\mathcal{L}_{\text{rank}}$ is ListMLE \cite{xia2008listwise} over the ground-truth permutation $\pi$,
\begin{equation}
\label{eq:listmle}
\mathcal{L}_{\text{rank}} = -\frac{1}{N}\sum_{i=1}^N \Big( \hat{s}_{\pi_i} - \log \sum_{j=i}^N \exp(\hat{s}_{\pi_j}) \Big),
\end{equation}
and $\mathcal{L}_{\text{pairwise}}$ is a margin-ranking loss ($\gamma = 0.05$) over pairs whose true impacts differ by more than $\gamma$. A general form with a maintainability term and a consistency term tying the heads to the explanation layer exists but is switched off throughout, so the learned engines and the explanation layer share no parameters (Supplementary Section~\ref{S-supp:objective}).

Models are trained with AdamW ($\eta = 3 \times 10^{-4}$, weight decay $10^{-4}$) under cosine warm restarts for up to 300 epochs, with early stopping at patience 30 on an inner validation split. Five seeds $\{42, 123, 456, 789, 2024\}$ are used throughout. Architectural hyperparameters follow conventional HGT values, and the loss coefficients were set by judgment. Neither was tuned on any evaluation split, so every learned arm is compared untuned.

\subsection{Ground-Truth Simulation Oracles}
\label{sec:4.3}

Ground truth comes from failure simulations over the raw structural multigraph $G_{\text{structural}}$. Table~\ref{tab:oracle_roles} summarizes the four oracles.

\begin{table}[htbp]
\centering
\footnotesize
\caption{Simulation oracles, operational constructs, and evaluation roles.}
\label{tab:oracle_roles}
\resizebox{\linewidth}{!}{%
\begin{tabular}{llll}
\toprule
\textbf{Oracle} & \textbf{Physical Mechanism} & \textbf{Nature} & \textbf{Role in Evaluation} \\
\midrule
$I^*(v)$ & BFS cascade reachability + QoS ladder & Seeded tie-breaking & Primary ranking target (RQ1--RQ3) \\
$I_{\text{comp}}(v)$ & Severity mixture: reachability + fragmentation & Deterministic & Explanation layer / Validate gate \\
$I_{\text{dyn}}(v)$ & Discrete-event SimPy message queuing & Stochastic & Convergent-validity probe \\
$I_M(v)$ & Reverse \texttt{DEPENDS\_ON} traversal & Deterministic & Unsupervised maintainability reference \\
\bottomrule
\end{tabular}%
}
\end{table}

\noindent\textbf{Primary target, $I^*(v)$.} The oracle crashes component $v$, propagates the outage through dependent topics, brokers and links by breadth-first traversal, and returns the mean fractional feed loss over the intact graph's subscriber population. A topic's feed loss is the fraction of its publishers that failed. It is scaled by a declared QoS severity ladder ($\times 1.2$ \texttt{RELIABLE}, $\times 1.15$ high priority, $\times 1.05$ medium) and clamped to $[0, 1]$. Five seeds break ties in propagation order, and $I^*(v)$ is their mean. It is reproducible from a fixed seed set, as CI gating requires.

\noindent\textbf{How much QoS is in this label.} The ladder reads reliability and priority only, but that does not bound the label's QoS content much: disabling QoS scaling entirely leaves the Application ordering nearly intact --- mean Spearman $\rho = 0.965$ against the ladder across the twelve folds (range $0.891$--$0.999$) --- and substituting a durability-aware $w(t)$ scaling moves it less still ($\rho = 0.977$). The top-$K$ set is the sensitive construct: ladder and topology-only labels agree at mean Jaccard $0.678$, so QoS changes \emph{which} components are named critical rather than their order. $I^*$ is therefore a near-topological target, which bounds what any QoS-encoding result can be credited with (Section~\ref{sec:rq2}).

\noindent\textbf{Further oracles.} $I_{\text{comp}}(v)$ is a severity-weighted mixture of reachability loss, fragmentation, throughput loss and flow disruption, with unswept AHP coefficients $(0.35, 0.25, 0.25, 0.15)$. It labels the explanation layer's evaluation and is never used for forecasting. $I_{\text{dyn}}(v)$ is a SimPy \cite{team2020simpy} message-flow simulation of emission rates, stochastic latencies and broker buffer saturation. It returns the drop in delivered message rate to surviving consumers and serves as an independent convergent-validity probe. It agrees with $I^*$ at $\rho = 0.627$, which is substantial but below $I^*$'s own seed-to-seed test--retest of $0.811$--$1.000$, so the two measure related but distinct constructs (Supplementary Section~\ref{S-supp:convergent}). $I_M(v)$ is a reverse-dependency traversal kept as a structural maintainability reference; it is never a training label. Topic criticality is a predictor input, so it is masked out of every oracle's severity term. Results established against one oracle are never transferred to another.

\subsection{Input--Label Independence Guarantee}
\label{sec:4.4}

Features are built only from $G_{\text{analysis}}$: static topology, code metrics and declared QoS. Labels are computed only on $G_{\text{structural}}$ by the simulation oracles. No simulation output or runtime telemetry is exposed as a predictor input, and a CI test enforces the separation (\path{tests/test_independence_guarantee.py}). This rules out circular feature construction. It does not make features and labels independent of the topology they both derive from: $I^*(v)$ is a reachability functional of the same architecture. The learning task is therefore to combine pre-computed structural cues into a ranking that matches the simulator, which is why the closed-form engine is a strong reference.

\section{The Explanation Layer: Standards-Grounded Criticality Attribution}
\label{sec:5}

A ranking says where risk is highest, not how to reduce it. A component may be critical because it is an unreplicated single point of failure, an error-propagating cascade hub, or a high-coupling maintainability bottleneck. Each of these calls for a different intervention: replication, circuit breakers, or decoupling. The explanation layer attributes these causes after ranking. It reads the same node properties (Section~\ref{sec:3.4}), shares no parameters with the engines, and is not used as a ranker; its ranking correlation is reported for reference in Supplementary Sections~\ref{S-supp:loso-gpu} and~\ref{S-supp:rq4}.

\subsection{Grounding in ISO/IEC Standards}
\label{sec:5.1}

Following ISO/IEC 25010:2023 \cite{iso25010_2023} and ISO/IEC 25019:2023 \cite{iso25019_2023}, criticality is profiled along \textbf{Reliability ($R$)}, split into \textbf{Fault Tolerance ($FT$)} and \textbf{Availability ($A$)}, and \textbf{Maintainability ($M$)}. $FT$ captures error-cascade potential and informs circuit breakers and redundancy. $A$ captures structural single points of failure and informs replication. $M$ captures coupling and code-level complexity and informs decoupling and refactoring. Safety and security, which need hazard logs, are out of scope.

\subsection{Composite Quality Score}
\label{sec:5.2}

Figure~\ref{fig:rm} summarizes the layer. All metrics are rank-normalized to $[0, 1]$ within the graph and combined with AHP-derived weights \cite{saaty1980ahp}:
\begin{itemize}\tightlist
    \item $FT(v) = 0.45 \cdot \text{RPR}(v) + 0.30 \cdot \text{Deg}_{\text{in}}(v) + 0.25 \cdot \text{CDPot}_{\text{enh}}(v)$, over Reverse PageRank, normalized in-degree and normalized cascade depth on $G_{\text{analysis}}^\top$;
    \item $A(v) = 0.25 \cdot \text{AP}_c^{\text{dir}}(v) + 0.20 \cdot \text{QSPOF}(v) + 0.20 \cdot \text{BR}(v) + 0.25 \cdot \text{CDI}(v) + 0.10 \cdot w(v)$, over directed articulation severity, QoS-weighted SPOF severity, bridge ratio, the Connectivity Degradation Index and the node QoS weight;
    \item $R(v) = r_{\text{FT}} \cdot FT(v) + (1 - r_{\text{FT}}) \cdot A(v)$ with $r_{\text{FT}} = 0.36$;
    \item $M(v) = 0.35 \cdot \text{BT}(v) + 0.30 \cdot w_{\text{out}}(v) + 0.15 \cdot \text{CQP}(v) + 0.12 \cdot \text{CouplingRisk}_{\text{enh}}(v) + 0.08 \cdot (1 - \text{CC}(v))$, over betweenness, QoS-weighted efferent coupling, code-quality penalty, coupling risk and clustering.
\end{itemize}
The composite is $Q(v) = 0.80 \cdot R(v) + 0.20 \cdot M(v)$, and an ISO/IEC 25019 context-of-use vector can reweight $R$ and $M$. Intra-dimension weights are shrunk towards a uniform prior ($\lambda = 0.70$). If $Q(v)$ is used to rank, a fully uniform prior is better ($0.319$ vs.\ $0.200$; Supplementary Section~\ref{S-supp:params}). The AHP matrices and their consistency diagnostics are in Supplementary Section~\ref{S-supp:ahp}. Components above the Tukey upper fence of $Q$ are flagged CRITICAL (mean $4.2\%$ of components). High $A$ with low $FT$ indicates a single point of failure that needs replication, while high $FT$ indicates a cascade hub that needs circuit breakers (example card: Supplementary Section~\ref{S-supp:card}).

\begin{figure}[htbp]
\centering
\includegraphics[width=\linewidth]{figures/Figure_4}
\caption{The explanation layer. Rank-normalized graph metrics feed the ISO/IEC 25010 sub-characteristics Fault Tolerance, Availability and Maintainability (CR: coupling risk; CC: clustering coefficient), which combine into Reliability and the composite $Q(v)$. A component above the Tukey fence of $Q$ is flagged, and its $FT$/$A$/$M$ profile names the remediation class.}
\label{fig:rm}
\end{figure}

\subsection{Counterfactual Verification of Remediation}
\label{sec:5.3}

Candidate repairs (broker replication, circuit-breaker insertion, topic decoupling) are generated for flagged components and verified counterfactually in memory. A repair is accepted only if it reduces systemic impact beyond seed noise ($\Delta \bar{I}_{\text{comp}} > \kappa \cdot \sigma_{\text{seed}}$, $\kappa \ge 1$) without introducing new articulation points. Developer studies and patch-synthesis benchmarks are future work.

\section{Experimental Setup}
\label{sec:6}

\subsection{Corpus and Replication Package}
\label{sec:6.1}

The corpus comprises 2,812 components across seventeen architectures (Table~\ref{tab:4}). Twelve synthetic topologies form the LOSO folds. They span autonomous vehicles, financial trading, healthcare, industrial SCADA, smart-city IoT, telecom RAN, logistics, gaming, microservices, enterprise integration and air-traffic management. All twelve come from one generator family, as do their code metrics, so LOSO measures transfer across configurations of that generator. Each regenerates byte-identically from a committed configuration, and CI verifies this against a SHA-256 manifest.

\begin{table}[htbp]
\centering
\small
\caption{Evaluation corpus. The twelve synthetic topologies are the LOSO folds; the five open-source system models are excluded from all training and used only for zero-shot transfer (Section~\ref{sec:rq3}). Counts are read from the committed topology files and verified in CI; per-scenario composition: Supplementary Section~\ref{S-supp:corpus}.}
\label{tab:4}
\begin{tabular}{lrrrrrrr}
\toprule
\textbf{Dataset} & \textbf{$|V|$} & \textbf{$|V_{\text{app}}|$} & \textbf{Topics} & \textbf{Brokers} & \textbf{Hosts} & \textbf{Libs} & \textbf{$|E|$} \\
\midrule
Synthetic evaluation scenarios (11) & 2{,}387 & 1{,}295 & 588 & 60 & 194 & 250 & 10{,}657 \\
Synthetic case study — ATM (1)    &    74 &    26 &  27 &  5 &   8 &   8 &   261 \\
\midrule
\textbf{Synthetic subtotal (12 LOSO folds)} & \textbf{2{,}461} & \textbf{1{,}321} & \textbf{615} & \textbf{65} & \textbf{202} & \textbf{258} & \textbf{10{,}918} \\
\textbf{Open-source system models (5)}    & \textbf{351}    &   \textbf{141} & \textbf{120} & \textbf{16} &  \textbf{32} &  \textbf{42} & \textbf{700} \\
\midrule
\textbf{Total} & \textbf{2{,}812} & \textbf{1{,}462} & \textbf{735} & \textbf{81} & \textbf{234} & \textbf{300} & \textbf{11{,}618} \\
\bottomrule
\end{tabular}
\end{table}

\noindent\textbf{The five open-source systems are hand-authored models.} Autoware.universe (ROS~2), EdgeX Foundry, Home Assistant, and two meshes modelled after Online Boutique and Train-Ticket were each written by one author as typed multigraphs from public documentation (\path{saag/adapters/realworld_adapter.py}). They are not mechanical extractions. Brokers, QoS profiles, code metrics and host specifications are partly assumed. Two models depart materially from their originals: the Online Boutique model is a 22-application pub-sub mesh with four brokers, whereas the original is about eleven gRPC services with no broker, and the Train-Ticket model represents its service-discovery server as a broker. Neither contains a synchronous call edge. RQ3 therefore tests transfer to independently authored architecture models, not to deployed systems.

\noindent\textbf{Replication package.} Datasets, harnesses, checkpoints and result artifacts are archived on Zenodo (see Data Availability). The public repository documents each experiment: its protocol, hyperparameters, \texttt{make} target, artifacts and the supplementary section holding its extended results (\sagexperiments).

\subsection{Predictors}
\label{sec:6.2}

SaG's closed-form engine (\texttt{Topo-QoS}), learned engines (\texttt{HGT-QoS}, \texttt{GAT-QoS}) and hybrid engines are compared against unweighted centrality (Topo) and untyped GNNs (Table~\ref{tab:predictor_taxonomy}). On a GNN, \texttt{-QoS} means the 16-D QoS edge vector (Section~\ref{sec:4.1.1}), and \texttt{Hybrid-X} is engine X corrected by the closed-form prior. \texttt{GAT} and \texttt{GAT-QoS} are untyped GATs matched to HGT in parameter budget, and \texttt{GAT-QoS} reads the same 16-D edge vector as \texttt{HGT-QoS}, so the four learned models form a $2\times2$ over typing (GAT vs.\ HGT) and the QoS channel. Smaller and projection-based variants that the study also ran are listed in Supplementary Section~\ref{S-supp:taxonomy}. All learned predictors ingest the native typed multigraph under LOSO. In QoS-off arms every edge weight is 1 and the QoS node columns are zeroed, though four centralities still carry QoS (Section~\ref{sec:3.4}).

\begin{table}[htbp]
\centering
\small
\caption{Predictors reported in this paper. \texttt{-QoS}: QoS-weighted distances (Topo) or the 16-D QoS edge vector (GNNs); \texttt{Hybrid-X}: engine X corrected by the \texttt{Topo-QoS} prior. Every predictor run in the study, with the earlier names used by the registered plan and the artifacts: Supplementary Section~\ref{S-supp:taxonomy}.}
\label{tab:predictor_taxonomy}
\resizebox{\linewidth}{!}{%
\begin{tabular}{llccccl}
\toprule
\textbf{Predictor} & \textbf{Evaluation Substrate} & \textbf{Typing} & \textbf{Edge Features} & \textbf{Parameters} & \textbf{Trained?} & \textbf{Empirical Role} \\
\midrule
\multicolumn{7}{l}{\textit{Training-free}} \\
\textbf{Topo} & $G_{\text{analysis}}$ (Flow Projection) & No & None & 0 & No & Standard centrality baseline \\
\textbf{Topo-QoS} & $G_{\text{analysis}}$ (Flow Projection) & No & Scalar $w(e)$ & 0 & No & SaG closed-form engine \\
\midrule
\multicolumn{7}{l}{\textit{Learned: typing $\times$ QoS channel at matched parameter budget}} \\
\textbf{GAT} & Native Multigraph & Homogeneous & None & 437{,}496 & Yes & Untyped, no QoS channel \\
\textbf{GAT-QoS} & Native Multigraph & Homogeneous & 16-D QoS Vector & 429{,}992 & Yes & Untyped SaG learned engine \\
\textbf{HGT} & Native Multigraph & Heterogeneous & Relation 1-hot; $w(e){=}1$ & 434{,}620 & Yes & Typed, no QoS channel \\
\textbf{HGT-QoS} & Native Multigraph & Heterogeneous & 16-D QoS Vector & 434{,}620 & Yes & Typed SaG learned engine \\
\midrule
\multicolumn{7}{l}{\textit{Hybrid engines: a learned engine corrected by the \texttt{Topo-QoS} prior}} \\
\textbf{Hybrid-HGT} & Native Multigraph & Heterogeneous & 16-D QoS Vector & 434{,}941 & Yes & \texttt{HGT-QoS} + prior \\
\textbf{Hybrid-GAT} & Native Multigraph & Homogeneous & 16-D QoS Vector & 431{,}433 & Yes & \texttt{GAT-QoS} + prior \\
\bottomrule
\end{tabular}%
}
\end{table}

\noindent\textbf{Closed-form scores.} Topo and \texttt{Topo-QoS} are computed on the Application--Library \texttt{DEPENDS\_ON} projection (Rules 1 and 5), because on the raw multigraph messages route through topics and brokers and Application betweenness vanishes:
\begin{equation}
\text{Topo}(v) = 0.6 \cdot \text{BT}(v) + 0.4 \cdot \text{AP}(v),
\qquad
\text{Topo-QoS}(v) = 0.6 \cdot \text{BT}_{w}(v) + 0.4 \cdot \text{AP}(v),
\end{equation}
where BT is normalized betweenness, $\text{BT}_{w}$ is betweenness over edge distances $d(e) = 1/(w(e) + 10^{-6})$, so strongly coupled edges attract shortest paths, and AP flags articulation points. In the evaluated implementation the articulation term reads zero for every node, so the reported scores rank exactly as betweenness and QoS-weighted betweenness. Restoring it lowers both (Supplementary Section~\ref{S-supp:ap-sensitivity}), so the evaluated form is the stronger reference and is kept as registered. The whole \texttt{Topo-QoS} gain over Topo therefore comes from QoS weighting of shortest paths. No predictor reads $G_{\text{structural}}$ (Section~\ref{sec:4.4}).

\subsection{Metrics, Protocols and Statistics}
\label{sec:6.3}

\noindent\textbf{Population.} Every predictor in a table is scored on the same node population, the Application set $V_{\text{app}}$. Pooling entity types conflates distinct base rates. Against $I_{\text{comp}}$, RM correlates at $\rho = 0.597$ on Applications but only $0.217$ pooled over all types (Supplementary Section~\ref{S-supp:detection}).

\noindent\textbf{Metrics.} Ranking is measured by Spearman $\rho$ against $I^*(v)$. Critical-set identification is measured by Overlap@$K$, the fraction of the true top-$K$ recovered by the predicted top-$K$ with $K = \text{round}(0.20\,|V_{\text{app}}|)$, at which top-$K$ precision, recall and $F_1$ coincide. PR-AUC, $F_1@\tau$ and nDCG@10 are reported in Supplementary Section~\ref{S-supp:identification}.

\noindent\textbf{Protocols.} Under \emph{LOSO}, models train on eleven scenarios and are tested zero-shot on the twelfth, over all 12 folds and five seeds, with equal 3-layer depth and inner-split early stopping. Under \emph{zero-shot transfer}, models trained on all twelve scenarios are evaluated without fine-tuning on the five system models. In-distribution node-split results are in Supplementary Section~\ref{S-supp:indist-table}.

\noindent\textbf{Statistics.} We use paired Wilcoxon signed-rank tests \cite{wilcoxon1945individual} and bootstrap 95\% CIs ($B = 2{,}000$) over folds \cite{efron1993introduction,spearman1904proof}; with 12 folds, the smallest attainable two-sided $p$ is $0.00049$. Folds share ten of eleven training scenarios, so the tests are anti-conservative \cite{nadeau2003inference}, and we read all $p$-values as nominal.

\noindent\textbf{Registered analysis plan.} Before the twelve-fold harness produced any result, we registered the primary contrast, \texttt{HGT-QoS} vs.\ \texttt{Topo-QoS}, together with \texttt{HGT} vs.\ \texttt{Topo-QoS} and Holm correction across the two. We call it \emph{registered} rather than pre-registered, because the plan is a file in our repository with no third-party timestamp. Three amendments registered further contrasts, each before its run and each Holm-corrected within its own family: the matched $2\times2$ (Amendment~2), Hybrid-HGT (Amendment~5) and Hybrid-GAT (Amendment~6). Every other contrast is exploratory. Because the sequence was adaptive (Amendment~6 followed Amendment~2's result), we also pool all eleven registered contrasts under one Holm correction (\path{reproduce/omnibus_holm.py}). Both hybrid primaries remain significant (Hybrid-GAT $p_{\text{omni}} = 0.016$, Hybrid-HGT $p_{\text{omni}} = 0.034$), and no other registered contrast reaches $\alpha = 0.05$ ($p_{\text{omni}} \ge 0.38$). Supplementary Section~\ref{S-supp:amendments} lists all six amendments with dates and outcomes.

\section{Results}
\label{sec:7}

All results are reported on the Application population ($V_{\text{app}}$) against the primary oracle $I^*(v)$, under the input--label independence guarantee (Section~\ref{sec:4.4}). Per-fold results, secondary strata and extended protocol notes are in the Supplementary Material and the experiment pages of the replication repository (Section~\ref{sec:6.1}). Figure~\ref{fig:results} summarizes the three main findings.

\begin{figure}[htbp]
\centering
\includegraphics[width=\linewidth]{figures/Figure_5}
\caption{Main results at a glance, Application population. (A)~Mean Spearman $\rho$ with 95\% bootstrap intervals under LOSO (filled circles; CPU sweeps of Table~\ref{tab:hybrid}) and zero-shot on the five system models (open diamonds). The hybrids lead on unseen synthetic architectures; the pure learned engines transfer best. (B)~Per held-out fold, the gain of \texttt{HGT-QoS} and of Hybrid-HGT over \texttt{Topo-QoS}; the arrow shows what the closed-form prior changes. It removes the learned engine's losses where the closed-form engine is strongest (Enterprise, Telecom RAN) and trims its largest gains where it is weakest. (C)~Cell means of the capacity- and channel-matched $2\times2$ (Table~\ref{tab:contrasts_matched}): the QoS channel raises both models by about $0.07$, while the typed and untyped lines stay together.}
\label{fig:results}
\end{figure}

\subsection{RQ1: SaG's Engines Against Structural Baselines}
\label{sec:rq1}
\label{sec:rq1-loso}

\paragraph{Summary} \emph{SaG's QoS-weighted closed-form engine (\texttt{Topo-QoS}, $\rho = 0.553$) outperforms unweighted centrality ($\rho = 0.349$) on all twelve held-out architectures ($+0.204$, $p = 0.0005$). On their own, the learned engines are statistically on par with it (\texttt{HGT-QoS} $0.622$, \texttt{GAT-QoS} $0.635$). The hybrid engines, in which a learned engine corrects the closed-form score, significantly outperform it ($+0.103$ and $+0.130$, each on 11/12 folds, Holm $p \le 0.0068$).}

Each of the twelve folds holds out one scenario and trains on the remaining eleven, and every predictor is scored on the same Application node set (26 to 300 nodes, $K$ between 5 and 60). Table~\ref{tab:hybrid} reports SaG's six engines and baselines from one family of CPU runs, in which the shared comparator rows are bit-identical.

\begin{table}[htbp]
\centering
\small
\caption{Main results. SaG's engines against the structural baselines under LOSO (twelve synthetic architectures, five seeds, Application population, CPU runs) and zero-shot on the five open-source system models (protocol of Table~\ref{tab:9b}). $\Delta\rho$ is paired by fold against the registered comparator \texttt{Topo-QoS}, with a bootstrap 95\% CI ($B = 2{,}000$) and a two-sided Wilcoxon test. $p_{\text{Holm}}$ is given for the hybrids, within each one's registered family (vs.\ \texttt{Topo-QoS} and vs.\ its own learned engine). The registered GPU sweep of \texttt{HGT-QoS}: Supplementary Section~\ref{S-supp:loso-gpu}. Per-fold values: Supplementary Section~\ref{S-supp:hybrid-folds}.}
\label{tab:hybrid}
\resizebox{\linewidth}{!}{%
\begin{tabular}{lccccccc}
\toprule
 & \multicolumn{5}{c}{\textbf{LOSO, twelve synthetic architectures}} & \multicolumn{2}{c}{\textbf{Five system models}} \\
\cmidrule(lr){2-6}\cmidrule(lr){7-8}
\textbf{Predictor} & \textbf{Mean $\rho$ [95\% CI]} & \textbf{$\Delta\rho$ vs \texttt{Topo-QoS} [95\% CI]} & \textbf{Won} & \textbf{$p$ ($p_{\text{Holm}}$)} & \textbf{Overlap@$K$} & \textbf{$\rho$ [95\% CI]} & \textbf{PR-AUC} \\
\midrule
\textbf{Topo} & 0.349 $[0.254, 0.452]$ & $-0.204$ $[-0.286, -0.122]$ & 0/12 & 0.0005 & 0.366 & 0.511 $[0.346, 0.703]$ & 0.474 \\
\textbf{Topo-QoS} & 0.553 $[0.443, 0.657]$ & --- & --- & --- & 0.388 & 0.526 $[0.357, 0.699]$ & 0.474 \\
\textbf{HGT-QoS} & 0.622 $[0.547, 0.690]$ & $+0.069$ $[-0.046, +0.174]$ & 8/12 & 0.266 & 0.426 & 0.760 $[0.714, 0.819]$ & 0.713 \\
\textbf{GAT-QoS} & 0.635 $[0.567, 0.696]$ & $+0.082$ $[-0.046, +0.201]$ & 7/12 & 0.233 & 0.438 & \textbf{0.805} $[0.759, 0.868]$ & \textbf{0.790} \\
\textbf{Hybrid-HGT} & 0.657 $[0.572, 0.733]$ & $+0.103$ $[+0.055, +0.152]$ & 11/12 & 0.0034 (0.0068) & 0.435 & 0.695 $[0.643, 0.730]$ & 0.602 \\
\textbf{Hybrid-GAT} & \textbf{0.683} $[0.603, 0.753]$ & $\mathbf{+0.130}$ $[+0.075, +0.190]$ & \textbf{11/12} & \textbf{0.0015} (\textbf{0.0029}) & \textbf{0.450} & 0.662 $[0.597, 0.727]$ & 0.600 \\
\bottomrule
\end{tabular}%
}
\end{table}

\noindent\textbf{The QoS-aware projection is the largest single gain.} Re-weighting shortest paths by declared QoS contracts lifts closed-form ranking from $\rho = 0.349$ to $0.553$ on every held-out architecture. Learned engines without the QoS channel reach only about this level (\texttt{HGT} $0.548$, \texttt{GAT} $0.563$; Section~\ref{sec:rq2}), so the representation carries much of the signal.

\noindent\textbf{On their own, learned engines are on par with the closed-form engine.} \texttt{HGT-QoS} ($+0.069$, $p = 0.266$) and \texttt{GAT-QoS} ($+0.082$, $p = 0.233$) lead \texttt{Topo-QoS} numerically, with intervals spanning zero. The registered confirmatory contrast, run on the GPU sweep the plan specified, agrees ($+0.085$, $p = 0.151$, Holm $0.303$; Supplementary Section~\ref{S-supp:loso-gpu}).

\noindent\textbf{Both hybrids significantly outperform the closed-form engine.} Hybrid-HGT reaches $\rho = 0.657$ ($+0.103$, Holm $p = 0.0068$) and Hybrid-GAT reaches $\rho = 0.683$ ($+0.130$, CI $[+0.075, +0.190]$, Holm $p = 0.0029$), each on 11 of 12 folds. Each meets the decision rule registered before its run. Both remain significant under one Holm correction pooled over all eleven registered contrasts of the study ($p_{\text{omni}} = 0.034$ and $0.016$; Section~\ref{sec:6.3}). They are the only engines in this study that significantly beat closed-form ranking, and Hybrid-GAT also has the highest Overlap@$K$ ($0.450$).

\noindent\textbf{Why the hybrids work: the engines are complementary.} \texttt{HGT-QoS} loses to \texttt{Topo-QoS} on four folds, all where the closed-form engine is strongest: Real-Time Gaming, Enterprise ($0.426$ vs.\ $0.795$), AV System and Telecom RAN ($0.427$ vs.\ $0.576$). Its largest gains come where the closed-form engine is weakest: Healthcare, IoT Smart City, ATM and Microservices, $+0.210$ to $+0.362$ (Figure~\ref{fig:results}B). The prior removes the failure mode: on Enterprise, Hybrid-HGT and Hybrid-GAT reach $0.735$ and $0.768$, and Telecom RAN turns from a loss into a win for both. The cost is a smaller gain on the weakest folds.

\noindent\textbf{Anchoring trades transfer for in-distribution accuracy.} On the five independently authored system models, both hybrids stay well above every training-free score ($0.695$ and $0.662$ vs.\ $0.511$--$0.526$) but below the pure learned engines ($0.760$ and $0.805$; Section~\ref{sec:rq3}). By the rule registered in Amendment~6, Hybrid-GAT therefore does not replace Hybrid-HGT as the recommended hybrid ($0.662 < 0.695$).

\noindent\textbf{Label noise and inert components.} Re-running the oracle across five seeds gives a test--retest rank correlation of $0.811$--$1.000$ (median $0.982$), well above every engine's mean $\rho$. Between $21\%$ and $52\%$ of each held-out population carries zero simulated impact. In the registered sweep, restricting evaluation to components with positive impact halves every predictor's correlation, learned or not, and leaves the method ordering unchanged (Supplementary Section~\ref{S-supp:loso-active}).

\subsubsection{How the Hybrids Are Built}
\label{sec:rq1-hybrid}

Each hybrid gives a learned engine the rank-normalized \texttt{Topo-QoS} score as one extra input per Application and Library, and adds a learned correction to that score on the logit scale, $\hat{I}^*(v) = \sigma\big(z(v) + \alpha\,\operatorname{logit}(p(v))\big)$, with one learnable $\alpha$ (Figure~\ref{fig:engines}a). Everything else matches the underlying engine. Hybrid-HGT is built on \texttt{HGT-QoS} (Amendment~5; $321$ extra parameters) and Hybrid-GAT on \texttt{GAT-QoS} (Amendment~6; $1{,}441$ extra parameters). Each was registered with its contrasts and decision rule before any run, with no setting tuned, and evaluated in its own CPU sweep with its comparators re-run in the same invocation.

\subsection{RQ2: What Learned Engines Need}
\label{sec:rq2}

\paragraph{Summary} \emph{With model capacity and edge-channel width matched, the QoS edge channel is what improves learned ranking ($+0.073$ main effect, 10 of 12 folds), and relation-specific weights add nothing beyond it (typing main effect $-0.014$, interaction $+0.001$). The untyped \texttt{GAT-QoS} ($\rho = 0.635$) performs as well as the typed \texttt{HGT-QoS} ($0.622$).}

A first comparison against small untyped GATs ($28{,}168$ parameters against HGT's $434{,}620$, reading at most a scalar edge weight) credited typing with a large gain ($+0.234$ without QoS). That gain is an effect of capacity and channel width (Supplementary Sections~\ref{S-supp:naive-2x2} and~\ref{S-supp:loso-gpu}), and the control registered in Amendment~2 removes both differences. \texttt{GAT} and \texttt{GAT-QoS} are untyped GATs at HGT's parameter budget ($437{,}496$ and $429{,}992$), and \texttt{GAT-QoS} reads the same 16-D edge vector as \texttt{HGT-QoS}, including the relation one-hot. All four matched arms ran in one CPU sweep, and the decision rule was fixed before any control result existed (Table~\ref{tab:contrasts_matched}).

\begin{table}[htbp]
\centering
\small
\caption{The $2\times2$ with capacity and edge-channel width matched (Amendment~2): \texttt{GAT} ($\neg$T$\neg$Q, $437{,}496$ parameters), \texttt{HGT} (T$\neg$Q, $434{,}620$), \texttt{GAT-QoS} ($\neg$TQ, $429{,}992$), \texttt{HGT-QoS} (TQ, $434{,}620$). One CPU sweep, twelve LOSO folds, five seeds, Application population. Holm correction across the three orthogonal quantities; simple effects are descriptive. Cell means: $0.563$, $0.548$, $0.635$, $0.622$.}
\label{tab:contrasts_matched}
\resizebox{\linewidth}{!}{%
\begin{tabular}{llrccccl}
\toprule
\textbf{Quantity} & \textbf{Contrast} & \textbf{$\Delta\rho$} & \textbf{95\% CI} & \textbf{Won} & \textbf{$W$} & \textbf{$p$} & \textbf{$p_{\text{Holm}}$} \\
\midrule
\multicolumn{8}{l}{\textit{Three orthogonal quantities, Holm-corrected across these three}} \\
\textbf{Typing (main effect)} & averaged over Q & $-0.014$ & $[-0.052, +0.023]$ & 4/12 & 29.0 & 0.470 & 0.940 \\
\textbf{QoS channel (main effect)} & averaged over T & $\mathbf{+0.073}$ & $[+0.013, +0.120]$ & 10/12 & 13.0 & 0.043 & 0.127 \\
\textbf{Typing $\times$ QoS interaction} & difference of differences & $+0.001$ & $[-0.050, +0.042]$ & 6/12 & 34.0 & 0.733 & 0.940 \\
\midrule
\multicolumn{8}{l}{\textit{Simple effects --- descriptive, not separately corrected}} \\
\textbf{Typing, QoS absent} & HGT vs.\ GAT & $-0.015$ & $[-0.064, +0.033]$ & 5/12 & 31.0 & 0.569 & --- \\
\textbf{Typing, QoS present} & HGT-QoS vs.\ GAT-QoS & $-0.013$ & $[-0.054, +0.026]$ & 4/12 & 27.0 & 0.380 & --- \\
\textbf{QoS channel, typing absent} & GAT-QoS vs.\ GAT & $\mathbf{+0.072}$ & $[+0.028, +0.109]$ & 10/12 & 9.0 & 0.016 & --- \\
\textbf{QoS channel, typing present} & HGT-QoS vs.\ HGT & $+0.073$ & $[-0.002, +0.136]$ & 10/12 & 19.0 & 0.129 & --- \\
\bottomrule
\end{tabular}%
}
\end{table}

\noindent\textbf{The QoS edge channel is the working ingredient.} At matched capacity, adding the 16-D QoS channel raises ranking by $+0.073$ with or without typing, on 10 of 12 folds each time, and significantly for the untyped pair ($+0.072$, CI $[+0.028, +0.109]$, $p = 0.016$). The channel also stabilizes training: the median within-fold seed spread of the untyped pair falls from $0.083$ to $0.010$ (the small GATs and HGT show the same pattern; Supplementary Section~\ref{S-supp:loso-gpu}).

\noindent\textbf{Relation-specific weights add nothing once capacity is matched.} \texttt{HGT} and \texttt{HGT-QoS} are within $0.015$ of their untyped counterparts \texttt{GAT} and \texttt{GAT-QoS} and win only 4--5 of 12 folds against them. Because \texttt{GAT-QoS} receives each edge's relation type as a feature, the precise finding is that relation-typed \emph{parameters} add nothing beyond relation-typed \emph{inputs}. Message directionality (the \texttt{HGT-QoS-U} control) remains unmatched.

\noindent\textbf{How much QoS the target can reward.} $I^*(v)$ is a near-topological target: a topology-only relabeling recovers its ordering at mean $\rho = 0.965$, and QoS acts mainly at its top-$K$ boundary (Section~\ref{sec:4.3}). On this target the QoS edge channel therefore acts largely as a relation-identity and coupling-strength signal. Oracles that express deadline misses, durability replay or priority inversion would let the encodings contribute contract semantics as well.

\subsection{RQ3: Zero-Shot Transfer to Models of Open-Source Systems}
\label{sec:rq3}

\paragraph{Summary} \emph{Learned engines trained only on synthetic scenarios transfer to independently authored system models: \texttt{HGT-QoS} reaches $\rho = 0.760$ $[0.714, 0.819]$ and its untyped counterpart \texttt{GAT-QoS} reaches $0.805$ $[0.759, 0.868]$, against $0.511$--$0.526$ for every training-free score, and both nearly double top-$K$ critical-set overlap. On components that actually propagate failures, the differences remain unresolved at five systems.}

The five systems are hand-authored models of Autoware.universe (ROS~2), EdgeX Foundry and Home Assistant, plus meshes modelled after Online Boutique and Train-Ticket (Section~\ref{sec:6.1}). They were written independently of the scenario generator but are models rather than extractions, and they carry labels from the same oracles. The test is therefore transfer to independently authored topologies under simulated reachability. \texttt{HGT-QoS} and \texttt{GAT-QoS} were trained on all twelve synthetic scenarios and evaluated zero-shot at the same 3-layer, 300-epoch budget as every other learned result. No system model contributed gradients or checkpoint selection. Where a system declares no QoS manifest, standard middleware defaults (ROS~2 Best-Effort/Reliable, MQTT QoS 0/1) are applied uniformly across all predictors.

\begin{table}[htbp]
\centering
\footnotesize
\setlength{\tabcolsep}{4.5pt}
\caption{Zero-shot transfer to hand-authored models of five open-source systems: Spearman $\rho$ against $I^*(v)$ on the Application population. Learned engines were trained on all twelve synthetic scenarios (3 layers, 300 epochs; five seeds, $\pm$ = spread over seeds). Training-free scores are deterministic and scored on identical labels and node sets. The last column restricts \texttt{HGT-QoS} to components with positive impact. All five models are encoded as publish--subscribe graphs; the grouping follows the paradigm of the \emph{original} system. Identification metrics: Supplementary Table~\ref{S-tab:identification}; bootstrap intervals and the active stratum for every predictor: Supplementary Section~\ref{S-supp:transfer-active}.}
\label{tab:9b}
\resizebox{\linewidth}{!}{%
\begin{tabular}{lrrrrccr}
\toprule
\textbf{System model} & \textbf{$|V_{\text{app}}|$} & \textbf{$n_{>0}$} & \textbf{Topo} & \textbf{Topo-QoS} & \textbf{HGT-QoS} & \textbf{GAT-QoS} & \textbf{HGT-QoS $\rho_{>0}$} \\
\midrule
\multicolumn{8}{l}{\textit{Originals are publish--subscribe systems}} \\
\textbf{Autoware.universe (ROS~2)} & 32 & 19 & 0.307 & 0.378 & 0.716 $\pm$ 0.081 & \textbf{0.758 $\pm$ 0.019} & $+$0.517 \\
\textbf{EdgeX Foundry (Industrial IoT)} & 22 & 10 & 0.534 & 0.534 & 0.793 $\pm$ 0.037 & \textbf{0.815 $\pm$ 0.055} & $+$0.183 \\
\textbf{Home Assistant (Smart Home)} & 24 & 17 & 0.297 & 0.289 & 0.864 $\pm$ 0.063 & \textbf{0.925 $\pm$ 0.023} & $+$0.702 \\
\midrule
\multicolumn{8}{l}{\textit{Originals are RPC systems (modelled as pub-sub)}} \\
\textbf{Online Boutique (pub-sub model)} & 22 & 8 & \textbf{0.891} & 0.888 & 0.710 $\pm$ 0.070 & 0.750 $\pm$ 0.119 & $-$0.031 \\
\textbf{Train-Ticket Booking Mesh} & 41 & 14 & 0.528 & 0.541 & 0.717 $\pm$ 0.096 & \textbf{0.777 $\pm$ 0.007} & $-$0.192 \\
\midrule
\textbf{Mean} & — & — & 0.511 & 0.526 & 0.760 & \textbf{0.805} & $+$0.236 \\
\bottomrule
\end{tabular}%
}
\end{table}

\noindent\textbf{Strong full-population transfer that does not depend on typing.} Both learned engines lead on 4 of 5 systems, with intervals that do not overlap those of the training-free scores (Table~\ref{tab:hybrid}). The untyped \texttt{GAT-QoS} beats \texttt{HGT-QoS} on all five systems ($\rho = 0.805$ vs.\ $0.760$). It also scores higher on Overlap@$K$ ($0.519$ vs.\ $0.470$) and PR-AUC ($0.790$ vs.\ $0.713$). Consistent with RQ2, what transfers is learning over SaG's QoS-annotated graph, not relation-specific parameters. Identification separates the engines most clearly: top-$K$ overlap averages $0.470$ against $0.248$ for the closed-form scores, and PR-AUC $0.713$ against $0.474$--$0.521$. On EdgeX, symmetric adapter-to-broker stars create betweenness ties that collapse closed-form triage entirely (Overlap@$K = 0.000$).

\noindent\textbf{Active components.} Restricted to components that propagate failures, \texttt{HGT-QoS} keeps a positive mean correlation ($\rho_{>0} = +0.236$) where every training-free score turns negative. Every interval spans zero at five systems, so this comparison is unresolved. $\rho_{>0}$ is positive on the three models of publish--subscribe systems and non-positive on the two modelled after RPC systems. Because both RPC-derived models are encoded as publish--subscribe graphs, this split cannot be attributed to call-tree semantics (Section~\ref{sec:limitations}). An earlier, non-blind 2-layer configuration leaves every conclusion unchanged (Supplementary Section~\ref{S-supp:taxonomy}).

\subsection{RQ4: Analysis Cost}
\label{sec:rq4}

\paragraph{Summary} \emph{Neural inference is effectively free ($56\,\text{ms}$ for a 2{,}000-node architecture, $0.02\%$ of pipeline time). Cost is dominated by deterministic feature extraction, specifically the $O(|V|^2 + |V||E|)$ Connectivity Degradation Index, and it tracks the density of the derived dependency projection. Cold extraction takes $2$--$18\times$ (median $5.6\times$) as long as the in-process cascade simulation, so SaG's advantage is avoiding staging infrastructure rather than saving CPU time.}

\begin{table}[htbp]
\centering
\footnotesize
\setlength{\tabcolsep}{3pt}
\caption{Per-stage latency of the inference pipeline across graph sizes (CPU, median of 3 runs; 5 for the forward pass). The analysis stage is stable across repeats (p10--p90 within $1\%$ of the median), while the forward pass is dominated by interpreter and dispatch overhead; the 249-node row still carries first-call warm-up.}
\label{tab:scale}
\begin{tabular}{ccccccc}
\toprule
\textbf{$|V|$} & \textbf{$|E|$} & \textbf{\shortstack{Analyze\\(s)}} & \textbf{\shortstack{Graph $\to$\\Tensor (s)}} & \textbf{\shortstack{HGT\\Forward (ms)}} & \textbf{\shortstack{Analyze :\\Forward}} & \textbf{\shortstack{Forward\\p10--p90 (ms)}} \\
\midrule
249   & 1{,}127  & 1.74   & 0.010 & 26.5 & 66× & 13.0--34.4 \\
499   & 2{,}402  & 8.32   & 0.022 & 16.4 & 509× & 15.6--16.4 \\
999   & 6{,}422  & 44.54  & 0.056 & 21.1 & 2{,}108× & 19.1--36.5 \\
1{,}998 & 19{,}301 & 239.34 & 0.157 & 56.2 & \textbf{4{,}259×} & 43.8--57.8 \\
\bottomrule
\end{tabular}
\end{table}

At 2,000 components the HGT forward pass takes $56\,\text{ms}$ against $239\,\text{s}$ for structural analysis (Table~\ref{tab:scale}). Because the analysis stage produces the node features the forward pass consumes, end-to-end evaluation of an unseen 2,000-component architecture takes about four minutes, of which the learned model is $0.02\%$. The $56\,\text{ms}$ is the marginal cost of re-scoring an already-analysed graph. The dominant term is the Connectivity Degradation Index, computed for every node in the main connected component. That is a correctness requirement: restricting CDI to articulation points leaves it identically zero in redundant multi-publisher topologies and makes $A(v)$ near-constant.

Across the corpus, the complete analysis gate (structural analysis plus 18 anti-pattern detectors) runs in $0.16$--$79.3\,\text{s}$ per scenario. That is $2.0$--$17.7\times$ (median $5.6\times$) the five-seed cascade labeling sweep measured in the same session, and the gate is more expensive on all twelve scenarios. The premium tracks the size of the derived projection rather than component count; Enterprise, with 300 applications on 120 topics, is the maximum (Supplementary Section~\ref{S-supp:gate-ratio}). On raw CPU time, direct simulation is therefore faster wherever its parameters are available. SaG additionally scores infrastructure components and dependency edges and avoids staging infrastructure. Training the four learned arms once took $7.7$ CPU-hours, amortized over every later evaluation.

\section{Discussion, Threats to Validity, and Limitations}
\label{sec:8}

\subsection{Practical Consequences}
\label{sec:8.1}

\noindent\textbf{The representation carries the signal.} The most robust result is that SaG's QoS-aware dependency projection makes criticality legible to simple and learned analyzers alike. Closed-form centrality on the projection improves every held-out architecture over its unweighted form ($+0.204$). For learned engines, the QoS edge channel is the component that matters ($+0.073$ at matched capacity). Practitioners therefore gain most from modeling their architecture with typed entities and declared QoS contracts, whichever engine they then run. On that representation, closed-form and learned engines fail on different architectures, so combining them is worth more than choosing between them.

\noindent\textbf{Choosing an engine.} Table~\ref{tab:guidance} summarizes where each instrument fits. The closed-form engine needs no training and is the natural default for lightweight CI gates. The hybrids are the most accurate choice for architectures resembling the training corpus: they keep the closed-form engine's strength on dense projections such as Enterprise while adding the learned engines' gains elsewhere. Hybrid-HGT is the registered recommendation because it transfers better of the two. For substantially different systems, a pure learned engine with the QoS channel transfers best; because relation-specific weights add nothing at matched capacity, the untyped \texttt{GAT-QoS} is the simpler choice. The explanation layer then names a remediation class for each flagged component.

\begin{table}[htbp]
\centering
\small
\caption{How the instruments in the SaG portfolio are best used, given the evidence in Section~\ref{sec:7}.}
\label{tab:guidance}
\resizebox{\linewidth}{!}{%
\begin{tabular}{@{}p{0.22\linewidth}p{0.32\linewidth}p{0.46\linewidth}@{}}
\toprule
\textbf{Instrument} & \textbf{Context} & \textbf{Role and evidence} \\
\midrule
\textbf{\texttt{Topo-QoS}} (Closed-form) & Lightweight CI gates & Training-free, $\rho = 0.553$ out of distribution; $+0.204$ over unweighted centrality on 12/12 folds. \\
\textbf{Learned + QoS channel} & Substantially different or irregular architectures & Best transfer (\texttt{GAT-QoS} $0.805$, \texttt{HGT-QoS} $0.760$) and identification (PR-AUC $0.71$--$0.79$). \\
\textbf{Hybrid-HGT / Hybrid-GAT} & Architectures resembling the training corpus & Best LOSO ranking ($\rho = 0.657$ / $0.683$); significantly above \texttt{Topo-QoS} ($+0.103$ / $+0.130$, 11/12 folds). \\
\textbf{RM explanation layer} & Refactoring and root-cause discussion & ISO/IEC 25010 attribution (Availability vs.\ Fault Tolerance vs.\ Maintainability). \\
\bottomrule
\end{tabular}%
}
\end{table}

\noindent\textbf{Where learned engines help.} On this corpus, the learned engines gain most on dense, irregular meshes (Microservices, ATM) and on symmetric stars that create betweenness ties for closed-form scores (EdgeX). They lose ground on the largest, densest projection (Enterprise), where three rounds of message passing cover less of the graph. Each of these observations rests on one to three folds or systems. The replication repository tabulates them as working hypotheses with candidate mechanisms.

\noindent\textbf{Computational sustainability.} Pre-deployment analysis avoids provisioning staging clusters for chaos sweeps \cite{calero2015green,verdecchia2023greenai}; we state this as infrastructure avoidance, not a measured energy saving. At base SoC power ($28\,\text{W}$), one pass of the analysis gate over all twelve scenarios costs at most $0.83\,\text{Wh}$, and training the four learned arms once costs about $0.22\,\text{kWh}$, both upper bounds from wall-clock time. Direct RAPL/NVML measurement \cite{lannelongue2021green,schmidt2021codecarbon} and incremental caching of structural metrics across commits are the main open levers.

\subsection{Threats to Validity}
\label{sec:threats}

\noindent\textbf{Construct validity.} All labels are simulator-derived rather than observed failures. The behavioral queue-flow oracle agrees with the primary cascade oracle at $\rho = 0.627$ (Section~\ref{sec:4.3}); much of that agreement concerns which components are harmless. Because $I^*(v)$ is a reachability functional of the topology the predictors read, a strong closed-form comparator is expected. Retargeting the LOSO contrasts on $I_{\text{dyn}}$ is the next experiment. The five system models are the most judgement-laden input: each was written by one author, no second modeler has re-derived them, and RQ3's figures hold only for these models. The replication package includes a re-modeling protocol and an agreement tool (\path{reproduce/model_agreement.py}) so that this check is cheap to run.

\noindent\textbf{Internal validity.} Predictors consume $G_{\text{analysis}}$ and oracles traverse $G_{\text{structural}}$, which is enforced in CI. Substrate, training set, depth and early stopping are matched across learned arms. All typing conclusions rest on the capacity- and channel-matched control. Message directionality remains unmatched, because the registered \texttt{HGT-QoS-U} control has not been run. No hyperparameter was tuned on an evaluation split.

\noindent\textbf{Robustness to free parameters.} Under Morris screening only two of ten declared constants, the AHP shrinkage $\lambda$ and $r_{\text{FT}}$, have appreciable influence ($\mu^* \approx 0.13$ against $\le 0.025$ for the rest), and no setting of the topic-weight or QoS sub-weight constants changes any reported comparison (Supplementary Sections~\ref{S-supp:params}--\ref{S-supp:ahp}).

\noindent\textbf{External validity.} The synthetic corpus comes from one generator family, and the system models are small (22--41 applications). Scaling beyond 2,000 nodes would benefit from incremental caching or mini-batching \cite{zeng2020graphsaint}.

\noindent\textbf{Conclusion validity and repeatability.} LOSO folds share training scenarios, so $p$-values are nominal \cite{nadeau2003inference} and are read alongside fold-level sign consistency and bootstrap intervals. At fixed code, seeds and device, every figure reproduces at its reported precision, and all training-free cells also reproduce across devices. Learned cells move across code revisions and devices (up to $0.172$ in a fold mean; \texttt{HGT-QoS} $0.041$), so learned figures are reported against the released artifacts, and every comparison is made within one sweep. The drift ledger is in the replication repository (\path{reproduce/rerun_drift.py}).

\subsection{Limitations and Future Work}
\label{sec:limitations}

The explanation layer's attributions have not been evaluated with developers or against injected faults, and no published learned-criticality model (FINDER \cite{fan2020keyplayers}, DrBC \cite{fan2019learning}) has been reproduced on this corpus. Next steps are: (1)~an independent re-model of at least two of the five systems, with inter-modeler agreement reported; (2)~combining the hybrids' in-distribution accuracy with the pure engines' transfer, for example by learning when to trust the prior, and running the directionality control; (3)~retargeting RQ1 and RQ2 on $I_{\text{dyn}}$; (4)~synchronous call edges and a backward-propagating oracle, so that RPC architectures can be modeled natively; (5)~extracting system models from deployment manifests; and (6)~validating rankings against production incident data.

\section{Conclusion}
\label{sec:9}

Declared QoS contracts are what make cascading-failure risk in publish--subscribe systems measurable before deployment. SaG turns them into a typed, QoS-weighted architecture graph, and both closed-form and learned engines rank failure impact better on it. On twelve held-out synthetic architectures, the representation alone raises closed-form ranking from $\rho = 0.349$ to $0.553$, on every fold. Closed-form and learned engines prove complementary, each strongest where the other is weakest. Hybrid engines, in which a learned model corrects the closed-form score, are the most accurate in the study ($\rho = 0.657$ and $0.683$). They outperform the closed-form engine on 11 of 12 folds under decision rules registered before their runs, and remain significant under a Holm correction over all eleven registered contrasts. Trained only on synthetic data, learned engines transfer zero-shot to independently authored models of five open-source systems ($\rho = 0.760$--$0.805$ against $0.511$--$0.526$) and nearly double top-$K$ critical-set overlap. A control matched in capacity and edge-channel width identifies the QoS edge encoding, not relation-specific parameters, as what learned engines need, so a simple untyped attention network suffices.

For practice, the most valuable step is to model the architecture with typed entities and declared QoS contracts. SaG then offers an engine for each context: the training-free closed-form engine for lightweight CI gates, a hybrid engine for architectures resembling the training corpus, and a QoS-aware learned engine for substantially different systems, with neural inference in milliseconds. Its ISO/IEC 25010 layer names the remediation each flagged component calls for. Because the synthetic corpus regenerates byte-identically and every reported value is reconciled against released artifacts, others can extend or challenge these results directly.

Three steps would widen their reach: learning when to trust the closed-form prior, so that one engine combines the hybrids' in-distribution accuracy with the pure engines' transfer; retargeting the engines on the behavioral queue-flow oracle and on independent re-models of the five systems; and extracting models from deployment manifests, so that rankings can be validated against production incidents.


\section*{Declarations}

\noindent\textbf{CRediT Authorship Contribution Statement.} \textbf{Ibrahim Onuralp Yigit:} Conceptualization, Methodology, Software, Validation, Formal analysis, Investigation, Data curation, Writing -- original draft, Visualization. \textbf{Feza Buzluca:} Conceptualization, Methodology, Validation, Writing -- review and editing, Supervision.
\textbf{Declaration of Competing Interest.} The authors declare no competing financial interests or personal relationships that could have influenced this work.
\textbf{Funding.} This research received no external grant.

\noindent\textbf{Data Availability.} The replication package (datasets, harnesses, checkpoints, scripts) is available on Zenodo under DOI \href{https://doi.org/10.5281/zenodo.14922108}{10.5281/zenodo.14922108} \cite{sag-replication-package}\nocite{khodabandeh2025linkpred} with \texttt{uv}/\texttt{pip} environments. The public repository documents every experiment reported here --- protocol, hyperparameters, reproduction command, artifacts and extended results --- at \sagexperiments. Synthetic datasets regenerate byte-identically. The deposit ships all artifacts backing reported tables as a dated bundle (\texttt{SaG\_JSS\_\allowbreak Results\_<stamp>}) with a \texttt{MANIFEST.json} recording SHA-256 digests, commit hashes, and corpus provenance. Four supplementary artifacts predate provenance stamping and carry no commit or corpus digest: \texttt{atm\_\allowbreak scale\_\allowbreak sweep\_\allowbreak v3.json} (S6), \texttt{qos\_\allowbreak label\_\allowbreak ablation.json} (Section~\ref{sec:4.3}), \texttt{threshold\_\allowbreak sensitivity\_\allowbreak v3.json} (S3) and \texttt{topic\_\allowbreak weight\_\allowbreak sensitivity\_\allowbreak v3.json} (S1); their correspondence to the corpus is asserted by the bundle rather than recorded in the file. The verification script (\texttt{reproduce/reconcile\_manuscript.py}) runs standalone against the deposit, mechanically verifying 511 reported figures in the manuscript and supplement against the JSON artifacts.

\section*{Declaration of Generative AI and AI-assisted technologies in the manuscript preparation process}

\noindent During the preparation of this work the authors used Anthropic's Claude to assist with typesetting, LaTeX formatting, and the development of analysis and reporting scripts in the replication package. After using this tool the authors reviewed and edited the content as needed and take full responsibility for the content of the published article. The study design, the choice of experiments, the interpretation of results, and all scientific claims are the authors' own.


\renewcommand{\bibfont}{\footnotesize\linespread{0.85}\selectfont}
\setlength{\bibsep}{0pt}
\bibliographystyle{elsarticle-num}
\bibliography{refs}

\end{document}
